\papersection{Generating candidate histories}\label{sec:generation}
The generator must extend each word by either letter and launch a scan of the word already present, all on the same orbit. Each launch must keep its word attached to its own scan. Otherwise a scan of one history could certify another.

\lead{Word generation}
Use $A=SS$ and $B=SA$ as letters. The generator $g$ appears at the end of every wrapped word. We set $\wrap{\varepsilon}=Sg$ and $\wrap{zw}=z\wrap{w}$ for $z\in\{A,B\}$. The seed $g\wrap{\sigma}$ stores the program and input in the suffix $\sigma$. Given $u=\wrap{w}$, the call $gu$ exposes three subtrees over successive rounds. The branches $R_Au$ and $R_Bu$ lead to new calls $g\wrap{Aw}$ and $g\wrap{Bw}$. The third branch, $Qu$, launches the scan of $w$. Induction on the number of prepended letters gives a launch for every word ending in $\sigma$. The appendix gives the terms and the reductions for one extension.

\lead{Packet recognition}
A \emph{packet} begins with the contraction
\[
\mathcopy{Q\unichar{10216}w\unichar{10217} \unichar{8594} (T\unichar{10216}w\unichar{10217})(L\unichar{10216}w\unichar{10217})}{%
 Q\wrap{w}\step(T\wrap{w})(L\wrap{w})}
\]
Its left child stores the word under the fixed marker $T$. Its right child runs the scan and always has a redex. Thus the packet stays together while the scan advances. The marker can first appear in this position only through a launch. More precisely, $(T\wrap{w})H$ occurs on the seed orbit exactly when $w$ ends in $\sigma$ and $H=\pi^a(L\wrap{w})$ for some age $a$. The appendix proves this from the shape of the marker. On an arbitrary tree the detector could accept the word of one packet paired with the scan of another. On the seed orbit, both children come from the same launch.
